% !TEX program = pdflatex
% THEME 1 — Facile — Corrigé
\documentclass[11pt,a4paper]{article}
\usepackage[utf8]{inputenc}
\usepackage[T1]{fontenc}
\usepackage{lmodern}
\usepackage[french]{babel}
\usepackage{amsmath,amssymb,mathtools}
\usepackage{icomma}
\usepackage{siunitx}
\sisetup{output-decimal-marker={,},per-mode=symbol}
\DeclareSIUnit{\atm}{atm}
\DeclareSIUnit{\bar}{bar}
\DeclareSIUnit{\mmHg}{mmHg}
\usepackage[version=4]{mhchem}
\usepackage{xcolor}
\usepackage{colortbl}
\usepackage[most]{tcolorbox}
\usepackage{enumitem}
\usepackage{array}
\usepackage{booktabs}
\usepackage{fancyhdr}
\usepackage[a4paper,margin=2.1cm,headheight=26pt,headsep=10pt]{geometry}

\definecolor{primary}{RGB}{150,98,162}
\definecolor{primarydark}{RGB}{92,52,110}
\definecolor{excol}{RGB}{0,135,90}
\definecolor{attncol}{RGB}{200,40,55}

\newtcolorbox{rep}[1][]{enhanced,breakable,colback=excol!5,colframe=excol,
  boxrule=0.8pt,arc=2pt,left=3mm,right=3mm,top=2mm,bottom=2mm,
  fonttitle=\bfseries,coltitle=white,title={#1}}

\pagestyle{fancy}\fancyhf{}
\fancyhead[L]{\small\textcolor{primarydark}{\textbf{Fiche 7} — Loi des gaz}}
\fancyhead[R]{\small\textcolor{primarydark}{\textsc{Thème 1 · Facile · Corrigé}}}
\fancyfoot[C]{\small\thepage}
\renewcommand{\headrulewidth}{0.8pt}
\renewcommand{\headrule}{\hbox to\headwidth{\color{primary}\leaders\hrule height \headrulewidth\hfill}}
\setlength{\parindent}{0pt}\setlength{\parskip}{3pt}

\newcounter{exo}
\newenvironment{cor}[1][]{\refstepcounter{exo}\medskip
  \noindent{\color{primarydark}\bfseries Corrigé \theexo.}\ \textit{#1}\par\nopagebreak\smallskip}{\par}
\newcommand{\rf}[1]{\textbf{\textcolor{attncol}{#1}}}

\begin{document}

\begin{center}
{\Large\bfseries\color{primarydark} Thème 1 — Niveau Facile — Corrigé}
\end{center}
\textcolor{primary}{\hrule height 1pt}
\medskip

\begin{cor}[Températures en kelvins]
On applique $T=\theta+273$ (et $\theta=T-273$).
\begin{itemize}[leftmargin=1.5em,itemsep=1pt]
  \item (a) $25+273=\rf{\qty{298}{\kelvin}}$ ;\quad (b) $100+273=\rf{\qty{373}{\kelvin}}$ ;\quad
        (c) $-40+273=\rf{\qty{233}{\kelvin}}$ ;\quad (d) $727+273=\rf{\qty{1000}{\kelvin}}$.
  \item (e) $350-273=\rf{\qty{77}{\degreeCelsius}}$.
  \item (f) $0-273=\rf{\qty{-273}{\degreeCelsius}}$ : c'est le \textbf{zéro absolu}, la température la plus
        basse possible (l'agitation des molécules y cesse). Aucune température ne peut lui être
        inférieure.
\end{itemize}
\end{cor}

\begin{cor}[Conversions de pression]
\begin{itemize}[leftmargin=1.5em,itemsep=2pt]
  \item (a) $\qty{2}{\atm}=2\times\qty{101325}{\pascal}=\rf{\qty{202650}{\pascal}}$ ;\quad
        en \si{\mmHg} : $2\times760=\rf{\qty{1520}{\mmHg}}$.
  \item (b) $\dfrac{380}{760}=\rf{\qty{0.50}{\atm}}$.
  \item (c) $\dfrac{202650}{101325}=\rf{\qty{2.0}{\atm}}$.
  \item (d) $\qty{1.5}{\bar}=1{,}5\times\qty{e5}{\pascal}=\qty{150000}{\pascal}=\rf{\qty{150}{\kilo\pascal}}$.
\end{itemize}
\end{cor}

\begin{cor}[Boyle-Mariotte]
La température ne change pas et $n$ est constant : c'est \textbf{Boyle-Mariotte}, $p_1V_1=p_2V_2$.
On isole $V_2$ :
\[ V_2=\frac{p_1V_1}{p_2}=\frac{\qty{1.0}{\atm}\times\qty{6.0}{\litre}}{\qty{3.0}{\atm}}
     =\rf{\qty{2.0}{\litre}}. \]
La pression a été multipliée par $3$, donc le volume a été \rf{divisé par $3$} ($6{,}0\to2{,}0$). Le
produit $pV=\qty{6.0}{\atm\litre}$ est conservé.
\end{cor}

\begin{cor}[Charles--Gay-Lussac]
Pression constante, $n$ constant, températures \emph{déjà en kelvins} : c'est \textbf{Charles},
$\dfrac{V_1}{T_1}=\dfrac{V_2}{T_2}$.
\[ V_2=V_1\frac{T_2}{T_1}=\qty{3.0}{\litre}\times\frac{\qty{450}{\kelvin}}{\qty{300}{\kelvin}}
     =3{,}0\times1{,}5=\rf{\qty{4.5}{\litre}}. \]
La température absolue a augmenté de $50\,\%$ ($300\to450$), donc le volume aussi ($3{,}0\to4{,}5$).
\end{cor}

\begin{cor}[Avogadro]
Pression et température fixées : le volume est proportionnel à la quantité de matière (\textbf{Avogadro},
$\dfrac{V_1}{n_1}=\dfrac{V_2}{n_2}$).
\[ V_2=V_1\frac{n_2}{n_1}=\qty{4.48}{\litre}\times\frac{\qty{0.50}{\mole}}{\qty{0.20}{\mole}}
     =4{,}48\times2{,}5=\rf{\qty{11.2}{\litre}}. \]
On a multiplié $n$ par $2{,}5$, donc le volume par $2{,}5$.
\end{cor}

\end{document}
